% Figure from MATH 590 notes, section 11. Compile with the notes preamble.
\begin{tikzpicture}[scale=1.2]
  \coordinate (x1) at (0,0); \coordinate (x2) at (1.9,0); \coordinate (y) at (1.02,0.25);
  \coordinate (e1) at (1.4569,0.3571); \coordinate (e2) at (0.6495,0.3553);
  \fill[figB!10] (x1) circle (1.5);
  \fill[figB!10] (x2) circle (1.3);
  \begin{scope} \clip (x1) circle (1.5); \fill[figB!22] (x2) circle (1.3); \end{scope}
  \draw[figB, thick, dashed] (x1) circle (1.5);
  \draw[figB, thick, dashed] (x2) circle (1.3);
  \fill[figR!18] (y) circle (0.3852);
  \draw[figR, thick, dashed] (y) circle (0.3852);
  \draw[black!55, thin] (x1) -- (y) (x2) -- (y);
  \draw[figR, thick] (y) -- (e1);
  \draw[figR, thick] (y) -- (e2);
  \fill[black] (x1) circle (1.2pt) node[below=1pt, font=\scriptsize] {$x_1$};
  \fill[black] (x2) circle (1.2pt) node[below=1pt, font=\scriptsize] {$x_2$};
  \fill[black] (y) circle (1.2pt);
  \node[font=\scriptsize] at (1.02,0.06) {$y$};
  \node[font=\scriptsize, figR] at (1.62,0.47) {$\delta_1$};
  \node[font=\scriptsize, figR] at (0.5,0.47) {$\delta_2$};
  \node[font=\scriptsize, figB, anchor=east] at (-1.1,1.25) {$B_d(x_1,\varepsilon_1)$};
  \node[font=\scriptsize, figB, anchor=west] at (2.85,1.1) {$B_d(x_2,\varepsilon_2)$};
  \node[font=\scriptsize, figR] at (1.1,-1.55) {$B_d(y,\delta)$};
  \draw[figR, thin] (1.1,-1.43) -- (1.04,-0.14);
\end{tikzpicture}
